\subsection{\ifzh 环境分层将报告回升定位于不同起点和变化幅度\else Environmental strata locate rebound within different baselines and changes\fi}
\ifzh
九个环境组合的标准化预测显示，报告概率的环境差异远大于同一组合内的时期变化（表~\ref{tab:envrebound}；图~\ref{fig:envrebound}）。低海拔低绿度组合的三个时期概率为 2.29\%、1.42\% 和 1.81\%，低海拔高绿度为 6.27\%、4.10\% 和 5.17\%；两者均表现为下降后部分回升。高海拔高绿度组合则为 42.48\%、38.01\% 和 44.98\%。在相同海拔层内，较高绿度对应更高报告基线；跨海拔比较时，山地组合从较高基线上维持较高报告水平。从两次飓风之间到 Fiona 后，低海拔低、高绿度组合分别增加 0.39 和 1.07 个百分点，高海拔高绿度组合增加 6.97 个百分点。因此，以绝对报告概率增量衡量，较明显的回升点估计出现在高海拔高绿度环境；低地高绿度环境虽有较高报告水平，其末期预测仍低于自身初期。后续分析据此检验山地高报告是否对应较高的占域持续概率，以及低地较绿地点是否在整体报告回升不足的背景下仍具有较高的占域持续概率。
\else
Standardized predictions across the nine environmental cells show larger environmental contrasts than within-cell period changes (Table~\ref{tab:envrebound}; Figure~\ref{fig:envrebound}). Low-elevation, low-greenness probabilities were 2.29\%, 1.42\% and 1.81\% across periods, versus 6.27\%, 4.10\% and 5.17\% at low elevation and high greenness; both partially rebounded after declining. The high-elevation, high-greenness cell gave 42.48\%, 38.01\% and 44.98\%. Higher greenness corresponded to a higher reporting baseline within elevation bands, while upland cells retained high reporting from an already high baseline. From Inter-Hurricanes to Post-Fiona, the low-elevation low- and high-greenness cells increased by 0.39 and 1.07 percentage points, respectively, compared with 6.97 points in the high-elevation, high-greenness cell. By absolute probability change, the larger rebound point estimate therefore occurs in high, green environments. Greener lowlands retain higher reporting levels but remain below their own initial prediction. The subsequent persistence analysis examines whether high upland reporting corresponds to higher occupancy persistence probabilities and whether greener lowlands also have higher persistence despite incomplete aggregate reporting rebound.
\fi
\inctab{environment_rebound}
\ifzh
空间聚类区间显示跨期增量的估计精度较低。低海拔高绿度组合的 Fiona 后相对初期差值为 $-1.10$ 个百分点，95\% 区间为 $[-2.62,0.42]$；高海拔高绿度为 $+2.50$ 个百分点，区间为 $[-5.09,10.09]$。高海拔低绿度仅有 144 条完整清单，而低海拔低绿度有 22,433 条，环境覆盖高度不均衡。当前数据因此更清楚地定位了高、低报告环境，对各环境层已经超越初期水平的估计则较宽。所有环境组合的末期相对初期区间均跨越零，削弱了“山地已恢复而低地尚未恢复”的确定性划分；较高报告基线也会影响绝对增量大小。监测应通过同一地点、相近努力的重复调查，同时追踪各层的绝对变化和相对基线变化，并补充山地低绿度等稀疏组合，检验当前梯度是否反映更广泛的地点变化。
\else
Spatially clustered intervals show lower precision for the temporal increments. The low-elevation, high-greenness Post-Fiona contrast was $-1.10$ percentage points (95\% interval $[-2.62,0.42]$), while the high-elevation, high-greenness contrast was $+2.50$ points ($[-5.09,10.09]$). Coverage was uneven. Only 144 complete checklists represented high elevation with low greenness, versus 22,433 at low elevation with low greenness. The data locate high- and low-reporting environments more precisely than they establish how far each has surpassed its initial level. All cell-specific Post-Fiona versus Pre-Mar\'ia intervals span zero, weakening a categorical claim that uplands have recovered while lowlands have not. Higher baselines also affect the size of absolute increments. Repeated surveys at the same sites with comparable effort should track both absolute and baseline-relative changes and improve coverage of sparse cells, including less-green uplands, to test how widely the fitted gradient extends across sites.
\fi
\begin{figure}[htbp]\centering
\incfig{fig26_environment_rebound}
\caption{\ifzh
环境组合的标准化报告概率及变化。前两图为时期概率，第三图为 Fiona 后减去 María 前的百分点差及空间聚类 95\% 区间。所有时期采用相同的层内协变量分布。
\else
Standardized reporting probabilities and changes by environmental cell. The first two panels show period probabilities; the third gives Post-Fiona minus Pre-Mar\'ia percentage points and spatially clustered 95\% intervals. All periods share the same within-cell covariate distribution.
\fi}\label{fig:envrebound}
\end{figure}
\FloatBarrier
